\begin{center}
\textsc{Gauss} an \textsc{Bessel}.\quad Göttingen, 15.~November 1822.
\end{center}

\begingroup
\makeatletter
\renewcommand{\@makefnmark}{\hbox{\normalfont\textsuperscript{*})}}
\long\def\@makefntext#1{\parindent1.5em\noindent\textsuperscript{*})\ #1}
\makeatother
\par\noindent\hspace*{\parindent}.\ .\ .\ .\ .\ .\quad Die außerordentlichen Schwierigkeiten, ein Dreiecksnetz durch die Lüneburger Heide zu führen, kannte ich schon aus \textsc{Epaillys} Bericht, der dies geradezu für unmöglich erklärt und seine Dreiecke, um den südlichen Teil von
\par\clearpage
\noindent352\hfill{\small\textbf{BRIEFWECHSEL MIT BESSEL.}}\par\medskip
\noindent Hannover mit Hamburg zu verbinden, über Bremen, die Weser hinunter und so die Elbe wieder hinauf geführt hat.
Und doch hatte er mir gegenüber große Vorteile. Er beobachtete durchweg oben in seinen Signaltürmen, wo er sich viel leichter eine weitere Aussicht verschaffen konnte als ich, der überall zu ebener Erde gemessen hat. Beim Freihauen in den Wäldern musste er viel weniger auf die Schonung der Bäume Rücksicht nehmen. Überall hatte er kostenlosen Transport, Arbeiter, die ihn nichts kosteten, usw., während ich alles ohne Ausnahme mit Bargeld bezahlen muss und oft sehe, wie die Kosten sehr, sehr weit über meine Erwartungen hinausgehen. Er hatte eine ganze Brigade von Ingenieuren als Gehilfen usw.
Diese Erwägungen ließen mich meine erste Erkundungsreise (gegen Ende April) mit einer gewissen Sorge antreten.
Anfangs ging es sogar besser als erwartet. Ich fand, dass Garssen und Falkenberg sich unmittelbar mit Deister und Lichtenberg verbinden ließen. Als ich aber untersuchte, wie sich die Dreiecke von diesen beiden Punkten aus weiter nach Norden bis Lüneburg führen ließen, erwies sich das Gelände als so widerspenstig, dass ich mehrmals an einem erfolgreichen Ausgang zweifelte und fürchtete, das ganze Unternehmen aufgeben zu müssen.
Das Land überall flach, keine beherrschenden Punkte, überall Wald, teils in großen Waldgebieten wie dem Hassel, dem Lüsing, dem Becklinger Holz usw., teils in unzähligen kleineren Kämpen, die sich schachbrettartig voreinanderschieben.
Die Versuche meines Gehilfen, auf der Westseite etwas Brauchbares zu finden, blieben ganz ohne Erfolg. Mir selbst gelang es nach den beschwerlichsten Versuchen schließlich, gleichsam im Herzen der Heide zwei Dreiecke festzulegen, 9.~10.~12 [Falkenberg-Hauselberg-Wulfsode] und 10.~12.~13 [Hauselberg-Wulfsode-Wilsede]. Zugleich wurde mir jedoch klar, dass ich \textls[120]{dieser} Art von Arbeit bald nicht mehr gewachsen sein würde. Daher verschob ich die weitere Suche nach einer Möglichkeit, diese Dreiecke mit den weiter südlich gelegenen zu verbinden und nach Norden fortzuführen, auf einen späteren Zeitpunkt, zu dem ich mehr Gehilfen (auf die erste Reise hatte ich nur einen Offizier mitgenommen) und alle meine Instrumente bei mir haben würde.
Ich begann daher Mitte Juni mit den eigentlichen Messungen in Lichtenberg, und meine Hoffnung hat sich später auch weitgehend erfüllt.
Da es sich als ganz unmöglich erwies, Garssen durch das Freihauen einer Schneise mit Hauselberg zu verbinden, wurde nach langem Suchen ein noch weiter nördlich gelegener Punkt 14 [Scharnhorst]
\par\clearpage
\noindent{\small\textbf{ZUR HANNOVERSCHEN TRIANGULATION.}}\hfill353\par\medskip
\noindent gefunden, der mehr Hoffnung bot.
Später zeigte sich jedoch, dass auch dies völlig unmöglich war, weil der Wald zwischen 10 [Hauselberg] und 14 [Scharnhorst] auf zu hohem oder vielmehr nicht tief genug gelegenem Gelände lag, und ich musste mich glücklich schätzen, noch einen anderen Punkt zu finden, Breithorn, mit dem es schließlich gelang.
Inzwischen war auch die Linie 9.~13 [Falkenberg-Wilsede] mit einer sehr großen Schneise geöffnet worden, und so hätte 10 [Hauselberg] ganz entfallen können, wenn man 11 [Breithorn] früher gekannt hätte.
Ich behielt jedoch den erstgenannten Punkt bei und halte es für einen überaus wertvollen Vorteil, dass in meinem System drei Vierecke vorkommen, in denen alle sechs Richtungen tatsächlich hin und zurück gemessen wurden.
Nachdem zuvor alle Winkel in jedem Dreieck auf die richtige Summe ausgeglichen worden waren, brauchte es nur noch sehr kleine Änderungen, meist unter \(0^{\prime\prime}{,}1\), um diese drei Vierecke vollständig in Übereinstimmung zu bringen.
Es wäre zu wünschen, dass man bei jeder Messung \textls[120]{solche} Prüfungen hätte.
Es gibt Messungen, bei denen die Summen der drei Winkel überall erstaunlich gut stimmen und eine solche Prüfung zeigt, dass manche Winkel um \(2^{\prime\prime}\) bis \(3^{\prime\prime}\) sicher falsch sind.
Tatsächlich kann jeder die Prüfung anhand der Winkelsumme vornehmen. Die Prüfung mit Hilfe der Diagonalen ist weniger leicht zugänglich, so einfach sie auch für einen Mathematiker ist. Man kann sich der Vermutung nicht erwehren, dass die erste Prüfung zuweilen dazu gedient haben mag, die Beobachtungen zwar nicht zu verfälschen, sie aber doch ein wenig \textls[120]{auszuwählen} (eine Neigung dazu bemerkt man selbst bei \textsc{Delambre}).
Nördlich von Wulfsode hat man zum Glück noch den Timpenberg gefunden, der sich unmittelbar mit Hamburg verbinden lässt. Der Versuch, eine Schneise von Timpenberg nach Lüneburg zu schlagen, ist gescheitert, weil das Gelände dazwischen nicht tief genug lag.
Die Richtungen von 10 nach 11 und von 11 nach 14 haben große Schneisen erfordert.
Ich machte es mir stets zur Regel, mit der Rechnung mit allen Messungen, so wie ich sie erhielt, Schritt zu halten (bis auf die allerletzte Zeile). Nur so war es möglich, alle Schneisen mit äußerster Genauigkeit anzulegen, sodass kein einziger Baumstamm unnötig gefällt wurde, oder möglichst früh mit Sicherheit zu erkennen, dass eine Schneise unmöglich war.
So kannte ich z.~B.\ den Tiefenwinkel genau im Voraus, unter dem 14 [Scharnhorst] in 10 [Hauselberg] oder 20 [Lüneburg in 19 [Timpenberg] erscheinen musste. Im ersten
\par\smallskip\noindent{\footnotesize IX.}\hfill{\footnotesize45}\par
\clearpage
\noindent354\hfill{\small\textbf{BRIEFWECHSEL MIT BESSEL.}}\par\medskip
\noindent Fall war schon das Gelände vor dem Wald zu hoch, im zweiten, nachdem eine schmale Schneise 2000 Schritt weit fortgeführt worden war.
Von Wilsede aus habe ich noch einen Punkt [Nindorf] 3000 Meter nördlich von Timpenberg festgelegt, der eine direkte Sichtverbindung zu 20 [Lüneburg], 21 [Lauenburg], 23 [Hamburg] hat und sich wahrscheinlich mit 19 [Timpenberg] verbinden lässt. Die ganze Strecke dorthin muss aber freigehauen werden.
Auf dem Timpenberg habe ich die Winkel zwischen 12 [Wulfsode], 13 [Wilsede], 23 [Hamburg] erst vorläufig gemessen. Alle weiter südlich gelegenen Punkte sind abgeschlossen, ebenso Wilsede.
In Scharnhorst bin ich zuletzt gewesen. Dieser Punkt ließe sich durch leicht anzulegende Schneisen unmittelbar mit Lichtenberg und Deister verbinden, wodurch Garssen entbehrlich würde. Es wäre jedoch unmöglich gewesen, die Eignung dieses Punktes zu Beginn festzustellen.
Dieser Punkt ist wie mehrere andere eine ganz unscheinbare Stelle, von der man gar nicht vermuten würde, dass sie so wertvoll ist.
In vielen Richtungen, z.~B.\ von Lichtenberg nach Garssen und von Lichtenberg nach Falkenberg, verläuft die Linie so knapp über den Hindernissen dazwischen, dass sich [der Zielpunkt] bei gewöhnlicher Refraktion nur \textls[120]{wenige Bogensekunden} darüber erhebt und nur zuweilen \(30^{\prime\prime}\) bis \(40^{\prime\prime}\) erreichte. Von Falkenberg nach Wulfsode und umgekehrt kam das Licht bei schwächerer Refraktion \textls[120]{gar nicht} herüber und stieg immer erst in den späteren Nachmittagsstunden so weit auf, dass man es sehen konnte.
Vom Lichtenberg selbst sah ich auf dem Falkenberg selten etwas und umgekehrt. Das Heliotroplicht schwebte fast immer am freien Himmel (Entfernung 85542 Meter). Trotz der großen Entfernung stimmen die Messungen ebenso gut oder vielmehr fast besser als bei ganz kleinen Entfernungen von 10000 Metern, bei denen das Heliotroplicht noch immer fast zu stark ist, obwohl der Spiegel bis auf eine Öffnung von wenigen Quadratmillimetern abgedeckt war.
Um Zeit zu sparen, habe ich häufig selbst auf Entfernungen von 3 bis 4 bis 5 Meilen meine \(3\frac{1}{2}\) Fuß hohen steinernen Sockel selbst anvisiert.
Dadurch sind die Messungen hin und wieder etwas weniger genau geworden, als wenn ich nur Heliotroplicht verwendet hätte, doch wäre ich dann in diesem Jahr bei Weitem nicht so weit gekommen.
Der größte Fehler in der Summe der drei Winkel betrug in diesem Jahr \(1^{\prime\prime}{,}76\), auch \textls[120]{überhaupt}\footnote{In 19 Dreiecken.} der größte nach dem
\par\clearpage
\noindent{\small\textbf{ZUR HANNOVERSCHEN TRIANGULATION.}}\hfill355\par\medskip
\noindent im Dreieck 4.~5.~6 [Hohehagen-Hils-Brocken], wo er \(3^{\prime\prime}{,}7\) beträgt und hauptsächlich dem schwierigen Anvisieren des Brockenhausturms zuzuschreiben ist, der sich bei Sonnenschein wegen seiner Phase nie gut mittig anvisieren lässt.
Gern würde ich die Winkel dieses Dreiecks noch einmal nachmessen, wenn die Zeit es erlaubte.
In diesem Jahr waren bei mir oft drei Heliotrope zugleich in Betrieb, von denen mein Sohn einen bediente. Es ist tatsächlich ein prachtvoller Anblick.
Die Mitteilung, die ein gewisser \textsc{Schubach} im [Astronomischen] Jahrbuch 1825 über die Heliotrope gegeben hat, beruht auf einem Irrtum und hat mit meinen Heliotropen gar nichts gemein. Diese sind kunstvoll gebaute Instrumente, deren Konstruktion mich erst \textls[120]{sehr viel} Mühe gekostet hat, die nun aber auch, wie ich glaube, nichts zu wünschen übrig lassen.
Ich habe von beiden Bauarten, die ganz verschieden voneinander sind, hier Exemplare für General \textsc{Müffling} anfertigen lassen. Auch \textsc{Gerling} hat zwei erhalten. Vielleicht kann ich bald in \textsc{Schumachers} Astronomischen Nachrichten Abbildungen davon zeigen.
Zur Signalübermittlung und um bei großen Entfernungen den gegenüberstehenden Heliotropen zunächst die Richtung zu zeigen, benutzte ich oft einen großen Spiegel von einem Fuß im Quadrat, der wiederum anders ausgerichtet wurde.
Nach der Beschreibung meiner Gehilfen war das ein höchst prachtvoller Anblick. Noch in vier Meilen Entfernung hat es dem bloßen Auge zuweilen \textls[120]{weh}getan, lange hinzusehen. --- Doch nun genug davon.

\par\noindent\hspace*{\parindent}Ihre Art, geodätische Beobachtungen zu behandeln, habe ich mit Vergnügen in \textsc{Schumachers} Astronomischen Nachrichten gesehen.
Sie wissen, dass mich dieser Gegenstand schon vor vielen Jahren beschäftigt hat.
Da sich Ihren Arbeiten kaum etwas hinzufügen lässt, hätte jene Veröffentlichung meine eigene Arbeit überflüssig gemacht, wenn unsere Wege sich gekreuzt hätten.
Die Art, wie ich diesen Gegenstand behandelt habe, ist jedoch von der Ihren grundverschieden, und so werde ich bei der künftigen Veröffentlichung meiner Messungen auch meine theoretischen Arbeiten ausführlich entwickeln.
Ich hoffe, darin manches Unerwartete bieten zu können.
Aber diese Untersuchungen hängen mit einem reichen, fast unerschöpflich reichen Gebiet zusammen, und bei dieser wie bei so vielen anderen Gelegenheiten empfinde ich oft mit tiefer Wehmut, wie meine äußeren Verhältnisse mich an weit ausgreifenden theoretischen Arbeiten hindern.
Wenn solche Arbeiten voll gedeihen sollen, muss man sich ihnen ganz widmen können und darf nicht
\par\smallskip\noindent\hfill{\footnotesize45\textsuperscript{*}}\par
\clearpage
\noindent356\hfill{\small\textbf{BRIEFWECHSEL MIT BESSEL.}}\par\medskip
\noindent stündlich durch so verschiedenartige Arbeiten wie das Halten von Vorlesungen und all die Kleinarbeit beim Beobachten und Berechnen der Beobachtungen usw.\ usw.\ daran gehindert werden.\ .\ .\ .\ .\ .\ .

\par\noindent\hspace*{\parindent}Die Anstrengungen im heißen Sommer haben mich oft äußerst mitgenommen, zuweilen so sehr, dass ich glaubte, ihnen zu erliegen.
Auch ist es eine große Belastung bei den Arbeiten in der öden Lüneburger Heide, dass man oft nur eine schlechte Unterkunft findet und selbst diese nur meilenweit vom Arbeitspunkt entfernt.
Bei kühlem Wetter, das mir körperlich besser bekommt, fühlte ich mich im Allgemeinen immer leidlich wohl, und jetzt kann ich über mein Befinden nicht klagen.\ .\ .\ .\ .\ .\ .

\par\noindent\hspace*{\parindent}Aus dem obigen Bericht können Sie selbst ersehen, was zur Vollendung meiner Triangulation noch fehlt. Einen eigentlichen Plan für den nächsten Sommer kann ich aber jetzt noch nicht machen. Vieles wird dabei auch von \textsc{Schumachers} Mitwirkung abhängen.
Es wäre \textls[120]{möglich}, dass ich, wenn ich mich auf das Unentbehrliche beschränke, schon im April und Mai die Dreiecke abschließe und dann hierher zurückkomme. Vielleicht messe ich im Juli hier und im August an einem weiter nördlich gelegenen Punkt (Celle, Harburg, Hamburg?) mit dem Zenitsektor.
Wenn ich aber die Kosten nicht scheuen muss, könnte ich auch noch den ganzen Sommer für die Triangulation aufwenden, um alles so vollkommen wie möglich zu machen.

\par\noindent\hspace*{\parindent}Vor einigen Tagen habe ich aus München ein Universalinstrument erhalten, das ich vor zwei Jahren vor allem zur Messung von Azimuten bestellt hatte.
Nach meiner vorläufigen Reduktion scheint sich inzwischen das aus meinem Passageninstrument abgeleitete und bis Hamburg übertragene Azimut um weniger als \(1^{\prime\prime}{,}5\) von dem zu unterscheiden, das \textsc{Schumacher} mir mitgeteilt hat (er hat auch vorläufig den Winkel zwischen Lüneburg und Wilsede gemessen, von wo ich ihm Heliotroplicht zusenden ließ). Eine Azimutmessung meinerseits an anderen Punkten dürfte damit ziemlich überflüssig sein.
Ob ich Anlass sehen werde, mit dem Universalinstrument auch Polhöhen an mehreren Punkten zu messen, wird sich zeigen, wenn ich das Instrument erst eine Zeit lang benutzt habe.
Es ist gut gearbeitet, scheint mir aber nicht besonders bequem zu benutzen zu sein, vor allem für ein kurzsichtiges Auge, das zum Ablesen nicht nah genug an die Stellkreise herankommt.
Auf jeden Fall muss es sehr ermüdend sein, viel damit zu beobachten.
Meine Hori-
\par\clearpage
\noindent{\small\textbf{ZUR HANNOVERSCHEN TRIANGULATION.}}\hfill357\par\medskip
\noindent zontalwinkel habe ich alle mit einem 12-zölligen Theodoliten gemessen, der ein ausgezeichnetes Instrument ist.
Mit dem Zenitsektor habe ich im letzten Winter nur wenige Messungen sozusagen zur Probe gemacht. In optischer Leistung und Bedienungskomfort steht er dem \textsc{Reichenbach}schen Meridiankreis weit nach.

\par\noindent\hspace*{\parindent}Meine Gradmessung als solche kann für sich genommen eigentlich kein sehr wichtiges Ergebnis liefern. Ich weiß aber nicht, wann mit der Vollendung der \textsc{Schumacher}schen zu rechnen ist.
Soweit ich weiß, hat er seine Dreiecke erst über etwa ein Drittel der ganzen Länge geführt. Auch die Zenitsektormessungen in Skagen wurden nur von seinem Gehilfen \textsc{Caroc} vorgenommen.
Die genaue Länge seiner Basis (rund 6000 Meter) kenne ich noch nicht. Erst in diesem Herbst wollte er sie mit Hamburg verbinden, in seinem letzten Brief erwähnt er aber gar nichts davon.
Ich könnte zwar selbst eine Basis messen, und die Linie von Breithorn bis Scharnhorst scheint auf ihrer ganzen Länge (11220 Meter) keine unüberwindlichen Hindernisse zu bieten.
Doch selbst abgesehen von den großen Kosten gestehe ich, dass ich vor einer so höchst langweiligen Arbeit zurückschrecke.
An meinen trigonometrischen Messungen konnte ich immer eine interessante Seite finden, da ihre tägliche Reduktion immer einige Unterhaltung bot\footnote{Meine Art zu beobachten bestand darin, immer das zu messen, was sich gerade gut messen ließ, ungeachtet dessen, ob es ein unmittelbarer Dreieckswinkel war. An manchen Stationen nahm ich einige Hilfspunkte hinzu, um auch dann nicht untätig zu sein, wenn nur ein einziger Heliotrop leuchtete. Die Franzosen haben, wie ich sehe, in Spanien etwas Ähnliches getan, die Messungen aber ganz falsch ausgeglichen.}.
Bei Gelegenheit habe ich außerdem alle sichtbaren Objekte mit sich schneidenden Visierlinien erfasst, um zur Geografie des Landes beizutragen. Ich muss sagen, dass mir diese Arbeit mit ihren täglichen Ausgleichungen so lieb wurde, dass das Bemerken eines neuen Kirchturms sowie das Ermitteln und Berechnen seiner Lage wohl ebenso viel Vergnügen machte wie das Beobachten eines neuen Gestirns.
(Vor Gott ist es am Ende wohl auch einerlei, ob wir die Lage eines Kirchturms auf einen Fuß oder die eines Sterns auf eine Bogensekunde genau bestimmt haben.)
Ist eine Basismessung aber erst einmal im Gang, so sehe ich gar nichts, was den Verstand reizen oder unterhalten könnte. Bei einer vielleicht zwei Monate dauernden anstrengenden täglichen Arbeit muss man sich allein mit dem Gedanken auf-
\par\clearpage
\noindent358\hfill{\small\textbf{BRIEFWECHSEL MIT BESSEL.}}\par\medskip
\noindent recht halten, dass es eben doch einmal geschehen \textls[120]{muss}, um am Ende eine Zahl zu haben.

\par\noindent\hspace*{\parindent}Die große Genauigkeit beim Messen horizontaler Winkel mit Heliotroplicht über ungeheure Entfernungen und die Genauigkeit, mit der sich absolute Azimute mit Universalinstrumenten messen lassen, lassen mich glauben, dass sich heutzutage in geeignetem Gelände mit großem Vorteil eine Längengradmessung durchführen ließe, bei der man den Längenunterschied nicht auf die Zeitbestimmung, sondern auf die Konvergenz der Meridiane stützt.
Von den Pyrenäen sieht man den Mont Ventoux, von dort die Alpen, von den Alpen bis in die Steiermark, von dort bis nach Ungarn usw.
Freilich wächst hierbei die Genauigkeit nicht wie beim Meridianbogen mit der gesamten Länge des Bogens (den eigentlich geodätischen Teil kann man dabei als fehlerfrei betrachten). Denkt man sich also Teilstücke von ungefähr gleicher Größe, so wächst die Genauigkeit beim Breitengrad wie die Anzahl der Teilstücke, beim Längengrad ebenfalls, wenn der Längenunterschied der Endpunkte astronomisch, z.~B.\ durch Sternbedeckungen, beobachtet wird. Bei der Konvergenz der Meridiane wächst sie hingegen nur wie die Quadratwurzel aus der Anzahl der Teilstücke (ungefähr). Doch glaube ich, dass dieses Verfahren bei \textls[120]{sehr großen} Teilstücken genauer ist, als man es z.~B.\ erreichen kann, wenn man den Längenunterschied mit Pulversignalen durch \textls[120]{Zeitbestimmung} zu ermitteln versucht.
Am vorteilhaftesten wäre dieses Verfahren in nördlichen Gegenden, wenn es dort sehr weite Sichtverbindungen gibt.
Sollten sich diese nicht in Norwegen und Schweden oder in einigen Gegenden Russlands finden lassen?

\par\noindent\hspace*{\parindent}.\ .\ .\ .\ .\ .\quad Ich bin noch beim Ausgleichen, was ich leider schon einmal ganz umsonst gemacht habe, weil von einer Station versehentlich eine fehlerhafte Tabelle übernommen worden war.
Bei meiner Vorgehensweise wirkt sich gewissermaßen jeder etwa in Wilsede gemessene Winkel auf alle anderen bis Göttingen hin aus.
Ich werde Ihnen künftig einige Ergebnisse mitteilen.
Nach vorläufiger Rechnung liegt Hamburg, wenn man für Göttingen \(51^{0}\,31^{\prime}\,48^{\prime\prime}{,}7\) annimmt, auf \(53^{0}\,33^{\prime}\,1^{\prime\prime}{,}76\) und \(0^{0}\,2^{\prime}\,2^{\prime\prime}{,}97\) östlich von Göttingen, wobei der absolute Wert vorläufig auf \textsc{Zachs} Basis beruht.
\par\noindent\hspace*{\parindent}.\ .\ .\ .\ .\ .
\endgroup
